%%%%%%%%%%%%%%%%%%%%%%%%%%%%%%%%%%%%%%%%%%%%%%%%%%%%%%%%%%%%%%%%%%%%%%%%%%%%%%%%
%%
\cleardoubleoddpage%  Make sure to start each chapter on a new odd page
\chapter{JAX Implementation and Efficiency Study}
\label{sec:jaximplementation}

\section{Why JAX for this Problem}

JAX \cite{bradbury2018jax} is a natural second framework for this
reproduction because the underlying learning problem -- a fixed-architecture
multilayer perceptron trained with a static loss and optimizer for a fixed
number of epochs on a fixed-size dataset -- is exactly the regime in which
JAX's just-in-time compilation and functional programming model are
expected to pay off: the training step can be compiled once via
\texttt{jax.jit} and then executed as compiled XLA code for the full epoch
loop, without the per-step Python dispatch overhead that a straightforward
PyTorch training loop incurs. The sweep in this project runs the training
loop on the order of $10^4$ times per framework across the full paper
matrix (8 cases $\times$ 6 architectures $\times$ 14 training sizes
$\times$ 12 trials, doubled for the NSB experiment's two targets $u$ and
$p$), so per-run compilation and dispatch overhead is not incidental -- it
is repeated at full sweep scale, making JAX's compilation model directly
relevant to overall wall-clock time rather than a micro-optimization.

\section{JAX Reimplementation Strategy}

The JAX implementation mirrors the PyTorch implementation's structure as
closely as the two frameworks' programming models allow: an equivalent
multilayer perceptron built from explicit parameter pytrees, a training
loop built around \texttt{optax.adam} combined with
\texttt{optax.exponential\_decay} for the learning-rate schedule, and the
same sweep driver contract (case, architecture, training size, trial seed)
as the PyTorch sweep, reading the identical \texttt{.npz} training/test
files so that both frameworks train on exactly the same data. The two
implementation decisions documented for PyTorch in
Section~\ref{sec:pytorchreimplementation} -- Keras-style
activation-agnostic He initialization and per-trial seeding via
\texttt{jax.random.PRNGKey(trial\_seed)} -- apply identically here; the
JAX default initializer (which otherwise resolves to a Xavier-family
initialization) and the earlier, pre-fix behavior of reusing a single
fixed seed across all twelve trials were both corrected to match the
paper's protocol before any sweep data was generated.

The paper's checkpoint/restore rule (Section~\ref{sec:methodology}) is
implemented with the same two triggers as the PyTorch version, but exploits
JAX's immutable pytrees: because JAX parameters cannot be mutated in
place, snapshotting the checkpoint is a plain reassignment of the
parameter pytree rather than an explicit deep copy, which is both simpler
and avoids any risk of the PyTorch-style aliasing bugs that an in-place
\texttt{state\_dict} copy could otherwise introduce.

\section{Efficiency Metrics}

The framework comparison collects, for every (case, architecture, training
size, trial) run in both frameworks: wall-clock training time, the final
and best training loss, and the sparse-grid test error in the
mass-matrix-weighted $Y$-norm (Section~\ref{sec:methodology}). These are
aggregated geometrically (Section~\ref{sec:methodology}) and compared
directly between frameworks in
Section~\ref{sec:experimentalresults}.

Measured over the full sweep (24{,}094 completed runs,
Section~\ref{sec:experimentalresults}), the training-time ratio between
frameworks is case-dependent rather than uniformly in JAX's favor. For
diffusion, PyTorch is consistently slower than JAX by 9--28\%, the ratio
growing with the parametric dimension $d$. For NSB, the picture is mixed:
JAX is still faster in most of the eight (case, target) combinations, but
PyTorch is marginally faster in two of them
(Table~\ref{tab:framework-ratio}). This is consistent with a JIT
compilation overhead effect rather than a raw execution-speed difference:
diffusion has a single output target and the four training runs per
(architecture, $m$) combination are the dominant cost, so \texttt{jax.jit}'s
one-time compilation cost per distinct input shape is amortized well. NSB
doubles the number of distinct (target, architecture) shape combinations
the sweep compiles against, so a larger share of each JAX case's total time
is compilation rather than compiled execution -- eroding, and in two cases
reversing, JAX's advantage. This project did not instrument peak memory
usage directly (no memory profiler was wired into the sweep driver), so no
quantitative memory comparison is reported; qualitatively, JAX's
ahead-of-time compilation model means its peak memory footprint per run is
reached earlier and held more predictably than PyTorch's eager execution,
but this was not measured.

\section{Experimental Setup for Fair Comparison}

To keep the comparison controlled, both frameworks are run against the
identical dataset files, the identical twelve per-case trial seeds, the
identical six architecture configurations, and the identical
checkpoint/restore and stopping rule -- the paper's full 60{,}000-epoch cap
and $5\times10^{-7}$ loss tolerance for diffusion, and the documented
reduced 20{,}000-epoch/$5\times10^{-6}$ budget for NSB
(Appendix~\ref{app:an-appendix}), applied identically to both frameworks
within each experiment so the PyTorch/JAX comparison itself is always
protocol-for-protocol even though the two experiments' protocols differ
from each other. Both sweeps are run sequentially on the
same machine rather than interleaved, so that neither framework benefits
from a warmer cache or an otherwise idle system at the other's expense.
No attempt is made to give JAX an advantage through additional
JAX-specific tuning (e.g.\ custom XLA flags) beyond what the paper's
training protocol already specifies, since the goal of the comparison is
protocol-for-protocol parity, not a best-case performance benchmark for
either framework individually.

\section{Observed Trade-offs: Pytorch vs JAX}

The two frameworks reach statistically indistinguishable final test errors
under the identical protocol: the geometric-mean PyTorch/JAX error ratio
lies between 0.986 and 1.006 across all twelve (case, target) combinations
in both experiments (Section~\ref{sec:experimentalresults},
Table~\ref{tab:framework-ratio}) -- well within a single trial's own
geometric standard deviation. This is exactly what identical
initialization, optimizer, schedule, and checkpoint rule predict, and it
is itself a useful correctness check: two independently implemented
training loops converging this closely on identical data is evidence that
neither implementation has a hidden bug that the other doesn't share.

The training-time trade-off is more interesting than a flat "JAX is
faster." JAX wins consistently for diffusion (9--28\% faster), where the
sweep's shape diversity is lower and per-run compilation cost is amortized
well. For NSB, where the two-target ($u$, $p$) structure doubles the
number of distinct compiled shapes, JAX's advantage shrinks and disappears
entirely in two of the eight cases -- PyTorch's eager execution has no
compilation cost to amortize in the first place, so it is less sensitive
to how many distinct architecture/shape combinations a sweep touches. The
practical implication for a project like this one is that JAX's
compilation model is a real, measurable advantage at this problem scale,
but the size of that advantage depends on how much shape diversity a given
sweep has -- a consideration that would matter for deciding which
framework to reach for on a new sweep design, not just a one-time
benchmark number.
%%
%%%%%%%%%%%%%%%%%%%%%%%%%%%%%%%%%%%%%%%%%%%%%%%%%%%%%%%%%%%%%%%%%%%%%%%%%%%%%%%%
